\subsection{核空间对偶上的测度}

设 F 是局部凸空间。令 \(T_{s}\) 为弱拓扑 \(\sigma(\mathbf{F}^{\prime},\mathbf{F})\)（定义在 \(F^{\prime}\) 上），令 \(T_{c}\) 为在 F 的紧凸子集上一致收敛的拓扑。根据 Mackey 定理（TVS, IV, §1, 第~1号, 定理 1），\(T_{s}\) 和 \(T_{c}\) 在 \(F^{\prime}\) 上与 F 和 \(F^{\prime}\) 之间的对偶性相容；因此，\(F^{\prime}\) 上每个介于 \(T_{s}\) 与 \(T_{c}\) 之间的局部凸拓扑 T 也具有同样性质。若 T 是这样的拓扑，并且 \(F^{\prime}_{T}\) 表示赋予 T 的 \(F^{\prime}\) 空间，我们将 F 视为 \(F^{\prime}_{T}\) 的对偶。因此，对于上述类型的所有拓扑 T，\(F^{\prime}\) 上的投影测度都相同；若 \(\mu\) 是这样的投影测度，则其傅里叶变换是 F 上的函数。

把 F 上由满足下述条件的连续半范数 N 定义的局部凸拓扑 S 称为萨佐诺夫拓扑：\(N^{2}\) 是 F 上的正二次型，并且存在 F 上的连续正二次型 H，使得 \(\operatorname{Tr}\left(\mathrm{N}^{2}/\mathrm{H}\right)<+\infty\)。拓扑 S 比 F 上给定的拓扑粗；若这两个拓扑相同，则称 F 为核空间。这类空间将在后文详细研究。

定理 2（Minlos）。--- 设 F 是局部凸空间，T 是 \(F'\) 上介于 \(T_{s}\) 与 \(T_{c}\) 之间的局部凸拓扑，\(\mu\) 是 \(F'_{T}\) 上的投影测度。假设  的傅里叶变换 \(\Phi\) 在 F 上关于萨佐诺夫拓扑连续，则 \(\mu\) 是 \(\mu\)\(F'_{T}\) 上的测度。

令 \(\varepsilon > 0\)。由于 \(\Phi\) 在 \(F\) 上关于萨佐诺夫拓扑连续，存在  上的两个连续正二次型 \(Q\) 和 \(H\)，使得 \(F\)\(\operatorname{Tr}(Q / H) < +\infty\)，并且

\[
\Phi (0) - \mathcal {R} \Phi (x) \leqslant \varepsilon / 6
\]

对于每个 \(x \in \mathbf{F}\)（满足 \(\mathbf{Q}(x) \leqslant 1\)），由第~8号命题 8，\(|\mathcal{R}\Phi(x)| \leqslant \Phi(0)\) 对所有 \(x \in \mathbf{F}\) 都成立，因而

\[
\Phi (0) - \mathcal {R} \Phi (x) \leqslant \varepsilon / 6 + 2 \Phi (0) \mathrm{Q} (x)\tag{63}
\]

对于所有 \(x \in \mathbf{F}\) 都成立。

令 \(r = \left(12\Phi(0)\mathrm{Tr}(\mathbf{Q}/\mathbf{H})\varepsilon^{-1}\right)^{1/2}\)，并以 K 表示由满足 \(K\)\(x' \in F_{\mathcal{T}}'\) 且 \(\langle x, x' \rangle^2 \leqslant r^2\mathrm{H}(x)\) 对所有 \(x \in F\) 都成立的元素  构成的集合。由于 \(H^{1/2}\) 是 F 上的连续半范数，集合 K 在 \(F\)\(K\)\(F_{\mathcal{T}}'\) 中等度连续且闭；因此根据 Ascoli 定理（GT, X, §2, 第~5号, 定理 2 的推论 1），它在 \(F_{\mathcal{T}}'\) 中是紧的。

设 V 是 \(F'_{T}\) 的闭线性子空间，且余维有限；则 V 是 F 中某个有限维线性子空间 T 的正交补。令 \(\mu_{V}\) 为 \(T'\) 上的测度，它是  上的投影测度 \(\mu\) 在映射 \(F'_{T}\) 下的像；这里 \(p_{V}\) 是 T 到 F 的典范嵌入的转置；其傅里叶变换是 \(\Phi\) 在 T 上的限制。最后，根据 Hahn--Banach 定理（TVS, II, §3, 第~2号, 定理 1 的推论 1），\(p_{\mathrm{V}}(\mathrm{K})\) 等于集合 \(C_{r}\)，其中 \(x' \in T'\) 由满足 \(\langle x, x' \rangle^{2} \leqslant r^{2}\mathrm{H}(x)\) 对所有 \(x \in T\) 都成立的  构成。由不等式 (63)，可将第~9号命题 10 应用于  上的测度 \(\mu_{V}\)：取 q 为 \(T'\)\(2\Phi(0)\mathrm{Q}\) 在 T 上的限制，取 h 为 H 在 T 上的限制。于是 \(\mathrm{Tr}(q/h) \leqslant 2\Phi(0)\operatorname{Tr}(\mathrm{Q}/\mathrm{H})\)，从而

\[
\mu_ {\mathrm{V}} \left(\mathrm{T} ^ {\prime} - \mathrm{C} _ {r}\right) \leqslant 3 \left(\frac {\varepsilon}{6} + 2 \Phi (0) \operatorname{Tr} (\mathrm{Q} / \mathrm{H}) r ^ {- 2}\right) = \varepsilon .
\]

由于 \(p_V\) 通过取商定义了从 \(F_{\mathcal{T}}' / V\) 到 \(T'\) 的同构，第~1号命题 1 于是表明，\(\mu\) 是 \(F_{\mathcal{T}}'\) 上的测度。

证毕

推论。--- 设 F 是桶状核空间，T 是 \(F'\) 上介于 \(T_{s}\) 与 \(T_{c}\) 之间的局部凸拓扑，\(\mu\) 是 \(F'_{T}\) 上的投影测度，\(\Phi\) 是 \(\mu\) 的傅里叶变换。要使 \(\mu\) 成为测度，\(\Phi\) 在 F 上连续是必要且充分的。

必要性由第~8号命题 9 得出，充分性由定理 2 得出。

备注。--- 设 F 是桶状空间，T 是 \(F'\) 上介于 \(T_s\) 与 \(T_c\) 之间的局部凸拓扑。对于 T，\(F'\) 的每个紧子集在较粗的拓扑 \(T_s\) 下都是紧的。反过来，设 K 是 \(F'\) 的一个在 \(T_s\) 下紧的子集。由于 F 是桶状的，K 是等度连续的（TVS, III, §4, 第~2号, 定理 1）；但根据 Ascoli 定理，\(F'\) 的每个等度连续子集在 \(T_c\) 下相对紧，因而在 T 下更是相对紧的，所以 K 包含于 \(F'\) 的一个在 T 下紧的子集。不难由此推出，从 \(F'_T\) 到 \(F'_{T_s}\) 的恒等映射在这两个空间上的测度集合之间定义了一个双射。

